\section{One symbol, one role within a declared scope}
\label{sec:latex-style-02}

Maintain a notation register during drafting.  For every symbol record its
standing, type, first definition, scope, and conflicts.  A glyph is not globally
reserved merely because one chapter assigns it a role.  For example,
\(\Sigma\) is a formula context in a displayed calculus and is also part of
\(\Sigma^0_n\) in the arithmetical hierarchy.  These roles are legitimate
because each scope is declared before use.  International conventions for
mathematical symbols are catalogued in \cite{ISO80000}; where the
registered notation of this document departs from them, the register
governs.
Within one declared scope, however, a glyph must not change standing silently.

Overloading \(\entails\) is controlled by argument shape:
\[
  \mathcal M,s\entails\varphi,
  \qquad
  \mathcal M\entails\sigma,
  \qquad
  \Gamma\entails\varphi.
\]
These respectively denote satisfaction, truth of a sentence in a structure,
and semantic consequence.  The propositional, modal, and matrix chapters
register the analogous shapes \(v\entails\varphi\),
\(\mathcal M\entails T\), \(\mathcal M,w\entails A\), and
\(\entails_{\mathbf K_3}\); no meaning outside these declared shapes is
assigned to the glyph.

The unification section of the automated-reasoning chapter
(\cref{sec:automated-reasoning-03}) uses
\(\theta_{\mathrm u},\rho_{\mathrm u},\eta_{\mathrm u}\) for substitutions.
The visible subscript distinguishes those maps from the formula \(\theta\) in
the Tseitin example while preserving the substitution definitions and
composition convention of the first-order resolution chapter.

